\begin{definitionbox}{Definition (Continuously Differentiable; Class \(C^1\))}
\begin{definition}[Continuously Differentiable; Class \(C^1\)]
\label{def:class-c1}
Let \(I\subseteq\mathbb{R}\) be an interval and let \(f:I\to\mathbb{R}\). The function \(f\) is \emph{continuously differentiable} on \(I\), written \(f\in C^1(I)\), if \(f\) is differentiable at every point of \(I\) and the derivative \(f':I\to\mathbb{R}\) is continuous on \(I\).
\end{definition}
\end{definitionbox}

\begin{remark*}[Standard quantified statement]
\[
  f\in C^1(I)
  \Longleftrightarrow
  \bigl(\forall x\in I:\operatorname{IsDifferentiable}(f,x)\bigr)
  \land
  \operatorname{IsContinuous}(f',I).
\]
\end{remark*}

\begin{remark*}[Interpretation]
The condition \(f\in C^1(I)\) is strictly stronger than differentiability. A function may be differentiable everywhere while its derivative fails to be continuous; Darboux's Theorem still constrains that failure, because derivatives cannot jump. Thus the gap between differentiability and \(C^1\) is oscillatory rather than jump-like.
\end{remark*}

\begin{dependencies}
\begin{itemize}
  \item \hyperref[def:derivative-at-a-point]{Definition: Derivative at a Point}
  \item \hyperref[def:continuous-at-point]{Definition: Continuity at a Point}
\end{itemize}
\end{dependencies}

\begin{exposition}
\phantomsection\label{ex:x2sin}
Let
\[
  f(x)=
  \begin{cases}
    x^2\sin(1/x), & x\neq 0,\\
    0, & x=0.
  \end{cases}
\]
At the origin,
\[
  f'(0)=\lim_{x\to0}\frac{x^2\sin(1/x)}{x}
  =
  \lim_{x\to0}x\sin(1/x)=0
\]
by squeezing. For \(x\neq0\),
\[
  f'(x)=2x\sin(1/x)-\cos(1/x).
\]
The first term tends to \(0\), but \(\cos(1/x)\) oscillates, so \(f'\) is
not continuous at \(0\). Thus \(f\) is differentiable on all of
\(\mathbb{R}\), but \(f\notin C^1(\mathbb{R})\). The failure is oscillatory,
not a jump, as Darboux's theorem requires.
\end{exposition}

\begin{definitionbox}{Definition (Class \(C^k\))}
\begin{definition}[Class \(C^k\)]
\label{def:class-ck}
Let \(I\subseteq\mathbb{R}\) be an interval, let \(k\in\mathbb{N}\), and let \(f:I\to\mathbb{R}\). The function \(f\) is of \emph{class \(C^k\)} on \(I\), written \(f\in C^k(I)\), if the derivatives \(f^{(0)},f^{(1)},\ldots,f^{(k)}\) exist on \(I\) and \(f^{(k)}\) is continuous on \(I\). By convention, \(C^0(I)\) is the class of continuous functions on \(I\).
\end{definition}
\end{definitionbox}

\begin{remark*}[Standard quantified statement]
\[
  f\in C^k(I)
  \Longleftrightarrow
  \bigl(\forall j\in\{0,1,\ldots,k\}: f^{(j)}\text{ exists on }I\bigr)
  \land
  \operatorname{IsContinuous}\bigl(f^{(k)},I\bigr).
\]
\end{remark*}

\begin{remark*}[Predicate reading]
\[
  \operatorname{IsClassC}(f,I,k).
\]
For a fixed interval \(I\), the classes are nested:
\[
  f\in C^{k+1}(I)\Longrightarrow f\in C^k(I).
\]
\end{remark*}

\begin{remark*}[Negated quantified statement]
\[
  f\notin C^k(I)
  \Longleftrightarrow
  \exists x\in I:\ \bigl[f^{(k)}\text{ fails to exist at }x\bigr]
  \vee
  \bigl[f^{(k)}\text{ exists on }I\land
  \neg\operatorname{IsContinuous}(f^{(k)},x)\bigr].
\]
\end{remark*}

\begin{remark*}[Negation predicate reading]
\(\neg\operatorname{IsClassC}(f,I,k)\): somewhere the \(k\)-th derivative
either does not exist or exists but is not continuous. A single point where
the top derivative is discontinuous certifies \(f\notin C^k(I)\).
\end{remark*}

\begin{remark*}[Interpretation]
The notation \(C^k(I)\) records \(k\) layers of differentiability together with continuity at the top layer. This is the regularity scale used to state how many times a theorem may differentiate a function without leaving the controlled setting.
\end{remark*}

\begin{dependencies}
\begin{itemize}
  \item \hyperref[def:higher-derivatives]{Definition: Higher Derivatives}
  \item \hyperref[def:class-c1]{Definition: Continuously Differentiable; Class \(C^1\)}
\end{itemize}
\end{dependencies}

\begin{definitionbox}{Definition (Smooth Function; Class \(C^\infty\))}
\begin{definition}[Smooth Function; Class \(C^\infty\)]
\label{def:class-cinfty}
Let \(I\subseteq\mathbb{R}\) be an interval and let \(f:I\to\mathbb{R}\). The function \(f\) is \emph{smooth}, or of \emph{class \(C^\infty\)}, written \(f\in C^\infty(I)\), if \(f\in C^k(I)\) for every \(k\in\mathbb{N}\). Equivalently,
\[
  C^\infty(I)=\bigcap_{k=0}^{\infty}C^k(I).
\]
\end{definition}
\end{definitionbox}

\begin{remark*}[Standard quantified statement]
\[
  f\in C^\infty(I)
  \Longleftrightarrow
  \forall k\in\mathbb{N}:\ f\in C^k(I).
\]
\end{remark*}

\begin{remark*}[Negated quantified statement]
\[
  f\notin C^\infty(I)
  \Longleftrightarrow
  \exists k\in\mathbb{N}:\ f\notin C^k(I).
\]
A single order at which differentiability or continuity breaks demotes the
function out of \(C^\infty\).
\end{remark*}

\begin{remark*}[Interpretation]
Smoothness is the regularity language imported by differential geometry. Smooth manifolds are assembled from \(C^\infty\) transition maps, and the tangent-map construction seeded by the differential uses this vocabulary.
\end{remark*}

\begin{dependencies}
\begin{itemize}
  \item \hyperref[def:class-ck]{Definition: Class \(C^k\)}
\end{itemize}
\end{dependencies}

\begin{definitionbox}{Definition (Real-Analytic Function; Class \(C^\omega\))}
\begin{definition}[Real-Analytic Function; Class \(C^\omega\)]
\label{def:class-comega}
Let \(I\subseteq\mathbb{R}\) be an interval and let \(f:I\to\mathbb{R}\). The function \(f\) is \emph{real-analytic}, or of \emph{class \(C^\omega\)}, written \(f\in C^\omega(I)\), if for every \(a\in I\) there exists \(r>0\) such that
\[
  f(x)=\sum_{n=0}^{\infty}\frac{f^{(n)}(a)}{n!}(x-a)^n
  \qquad\text{for every }x\in(a-r,a+r)\cap I.
\]
\end{definition}
\end{definitionbox}

\begin{remark*}[Standard quantified statement]
\[
  f\in C^\omega(I)
  \Longleftrightarrow
  \forall a\in I\;\exists r>0\;\forall x\in(a-r,a+r)\cap I:\;
  f(x)=\sum_{n=0}^{\infty}\frac{f^{(n)}(a)}{n!}(x-a)^n.
\]
\end{remark*}

\begin{remark*}[Predicate reading]
The displayed quantified statement is the predicate-level form of $label: its hypotheses name the input conditions, and its conclusion names the asserted structure or relation.
\end{remark*}

\begin{remark*}[Interpretation]
Analytic regularity means that the Taylor series does more than approximate the function: near each point, the series equals the function. The infinite Taylor series and its convergence are developed in \hyperref[def:taylor-polynomial-at-a-point]{Taylor's construction and error analysis}, which supplies the series machinery behind this definition.
\end{remark*}

\begin{dependencies}
\begin{itemize}
  \item \hyperref[def:class-cinfty]{Definition: Smooth Function; Class \(C^\infty\)}
  \item \hyperref[def:taylor-polynomial-at-a-point]{Definition: Taylor Polynomial at a Point}
\end{itemize}
\end{dependencies}

\begin{theorembox}{Theorem (The Smoothness Tower)}
\begin{theorem}[The Smoothness Tower]
\label{thm:smoothness-tower}
For every nondegenerate interval \(I\), the inclusions
\[
  C^0(I)\supseteq C^1(I)\supseteq C^2(I)\supseteq \cdots
  \supseteq C^\infty(I)\supseteq C^\omega(I)
\]
hold, and each displayed inclusion is strict.
\hyperref[prf:smoothness-tower]{Go to proof.}
\end{theorem}
\end{theorembox}

\begin{remark*}[Standard quantified statement]
\[
  \forall k\in\mathbb{N}:\ C^{k+1}(I)\subsetneq C^k(I),
  \qquad
  C^\omega(I)\subsetneq C^\infty(I).
\]
\end{remark*}

\begin{remark*}[Failure modes]
\begin{description}
  \item[Exposition.]
The inclusions do not fail; only their strictness needs witnesses.
\begin{itemize}
  \item \(C^k\setminus C^{k+1}\): \(\varphi_k(x)=x^k|x|\), whose \(k\)-th
        derivative is \((k+1)!\,|x|\), continuous but not differentiable at
        \(0\).
  \item \(C^\infty\setminus C^\omega\): the flat function
        \(e^{-1/x^2}\), smooth with all derivatives zero at \(0\), so its
        Maclaurin series is zero while the function is not.
\end{itemize}
These separators prove every displayed containment is a genuine drop.
\end{description}
\end{remark*}

\begin{remark*}[Interpretation]
The inclusions are definitional; strictness is witnessed by examples. On any nondegenerate interval, translated and rescaled copies of \(x\mapsto x^k|x|\) separate \(C^k\) from \(C^{k+1}\). Translated copies of the flat function
\[
  f(x)=
  \begin{cases}
    e^{-1/x^2}, & x\neq 0,\\
    0, & x=0,
  \end{cases}
\]
separate \(C^\infty\) from \(C^\omega\): all derivatives at the flat point vanish, so the Taylor series is identically zero, while the function is positive on one side of that point.


Between \(C^1\) and \(C^2\) there is a regularity floor worth naming. A
\(C^1\) function has a continuous derivative; a \(C^2\) function has a
derivative that is itself differentiable. Many functions used in analysis
live in the gap: their first derivative is continuous and controlled, but a
classical second derivative may fail to exist at a few points. The control is
a Lipschitz bound on \(f'\), and the class is written \(C^{1,1}\).
\end{remark*}

\begin{dependencies}
\begin{itemize}
  \item \hyperref[def:class-ck]{Definition: Class \(C^k\)}
  \item \hyperref[def:class-cinfty]{Definition: Smooth Function; Class \(C^\infty\)}
  \item \hyperref[def:class-comega]{Definition: Real-Analytic Function; Class \(C^\omega\)}
  \item \hyperref[thm:darboux]{Theorem: Darboux's Theorem}
\end{itemize}
\end{dependencies}

\begin{figure}[H]
\centering
\resizebox{0.76\linewidth}{!}{%
\begin{tikzpicture}[scale=0.82]
  \def\lf{2.0+1.05*sin(80*\x)*0.9+0.16*\x}
  \glowcurve{color=atlasblue, domain=0.0:4.2, axis=0, top=4.2, expr={\lf}}
  \draw[atlas axes] (0,0)--(4.8,0);
  \draw[atlas axes] (0,0)--(0,4.4);
  \def\L{0.85}
  \draw[atlas probe,atlasconegreen] (0,3.6)--(4.2,{3.6-\L*4.2});
  \draw[atlas probe,atlasconegreen] (0,0.5)--(4.2,{0.5+\L*4.2});
  \draw[atlas probe,atlasconegreen] (0,1.0)--(4.2,{1.0+\L*4.2});
  \draw[atlas probe,atlasconegreen] (0,3.1)--(4.2,{3.1-\L*4.2});
  \node[atlas label,above] at (2.1,4.15){\scriptsize$|f'(x)-f'(y)|\le L|x-y|$};
  \node[atlas label,left] at (0,3.9){\scriptsize$-L$};
  \node[atlas label,left] at (0,0.5){\scriptsize$\pm L$};
  \node[atlas label,right] at (4.2,4.0){\scriptsize$\pm L$};
  \node[atlas label,right] at (4.2,0.6){\scriptsize$\pm L$};
  \node[atlas label,right] at (3.4,3.1){\scriptsize$y=f(x)$};
  \node[atlas label] at (2.5,1.0){\scriptsize$-L$};
\end{tikzpicture}%
}
\caption{A Lipschitz derivative changes at a uniformly bounded rate.}
\label{fig:lipschitz-derivative}
\end{figure}


\begin{definitionbox}{Definition (Lipschitz-Derivative Class \(C^{1,1}\))}
\begin{definition}[Lipschitz-Derivative Class \(C^{1,1}\)]
\label{def:class-c11}
Let \(I\subseteq\mathbb{R}\) be an interval and let \(f:I\to\mathbb{R}\).
The function \(f\) is of \emph{class \(C^{1,1}\)} on \(I\), written
\(f\in C^{1,1}(I)\), if \(f\in C^1(I)\) and \(f'\) is Lipschitz on \(I\):
there exists \(L\geq 0\) such that
\[
  |f'(x)-f'(y)|\leq L|x-y|
  \qquad\text{for all }x,y\in I.
\]
\end{definition}
\end{definitionbox}

\begin{remark*}[Standard quantified statement]
\[
  f\in C^{1,1}(I)
  \Longleftrightarrow
  f\in C^1(I)\land
  \exists L\geq 0\;\forall x,y\in I:\ |f'(x)-f'(y)|\leq L|x-y|.
\]
\end{remark*}

\begin{remark*}[Predicate reading]
\[
  f\in C^{1,1}(I)
  \Longleftrightarrow
  \operatorname{IsClassC}(f,I,1)\land
  \exists L\geq0:\operatorname{IsLipschitz}(f',I,L).
\]
The smallest admissible \(L\) is the Lipschitz constant of \(f'\):
\[
  \sup_{x\neq y}\frac{|f'(x)-f'(y)|}{|x-y|}.
\]
\end{remark*}

\begin{remark*}[Negated quantified statement]
\[
  f\notin C^{1,1}(I)
  \Longleftrightarrow
  f\notin C^1(I)\vee
  \bigl(\forall L\geq0\;\exists x,y\in I:\ |f'(x)-f'(y)|>L|x-y|\bigr).
\]
\end{remark*}

\begin{remark*}[Interpretation]
\(C^{1,1}\) says that \(f'\) behaves as if \(|f''|\leq L\), without requiring
\(f''\) to exist everywhere. It is the exact regularity tracked by many
first-order estimates: the derivative is not merely continuous, but changes
at a bounded rate.
\end{remark*}

\begin{dependencies}
\begin{itemize}
  \item \hyperref[def:class-c1]{Definition: Continuously Differentiable; Class \(C^1\)}
  \item \hyperref[cor:derivative-bound-implies-lipschitz]{Corollary: Derivative Bound Implies Lipschitz}
\end{itemize}
\end{dependencies}

\begin{theorembox}{Theorem (Placement of \(C^{1,1}\))}
\begin{theorem}[Placement of \(C^{1,1}\)]
\label{thm:c11-placement}
For every nondegenerate interval \(I\),
\[
  C^{1,1}(I)\subsetneq C^1(I),
\]
and \(C^{1,1}(I)\) is not contained in \(C^2(I)\). If \(K\) is a compact
interval, then
\[
  C^2(K)\subseteq C^{1,1}(K)\subsetneq C^1(K),
\]
and the first inclusion is strict.
\hyperref[prf:c11-placement]{Go to proof.}
\end{theorem}
\end{theorembox}

\begin{remark*}[Standard quantified statement]
\[
  \bigl(f\in C^{1,1}(I)\Rightarrow f\in C^1(I)\bigr),
  \qquad
  \bigl(K\text{ compact}\land f\in C^2(K)\Rightarrow f\in C^{1,1}(K)\bigr),
\]
with the displayed containments strict.
\end{remark*}

\begin{remark*}[Failure modes]
\begin{description}
  \item[Exposition.]
The global inclusion \(C^2(I)\subseteq C^{1,1}(I)\) is false on arbitrary
intervals: \(f(x)=x^3\) is \(C^2(\mathbb{R})\), but
\(f'(x)=3x^2\) is not Lipschitz on \(\mathbb{R}\). Strictness is witnessed by
two elementary functions.
\begin{itemize}
  \item \(C^{1,1}\setminus C^2\): \(f(x)=x|x|\) has \(f'(x)=2|x|\), which is
        Lipschitz, but \(f''(0)\) does not exist.
  \item \(C^1\setminus C^{1,1}\): \(f(x)=|x|^{4/3}\) has continuous first
        derivative, but \(f'\) has unbounded difference quotients at \(0\).
\end{itemize}
\end{description}
\end{remark*}

\begin{remark*}[Interpretation]
\(C^{1,1}\) is a genuine intermediate regularity condition, but it is not
the same as \(C^2\). A bounded second derivative forces the class; a
Lipschitz first derivative may have corners; and a continuous first
derivative may still vary too sharply to be Lipschitz.
\end{remark*}

\begin{dependencies}
\begin{itemize}
  \item \hyperref[def:class-c11]{Definition: Lipschitz-Derivative Class \(C^{1,1}\)}
  \item \hyperref[thm:smoothness-tower]{Theorem: The Smoothness Tower}
  \item \hyperref[cor:derivative-bound-implies-lipschitz]{Corollary: Derivative Bound Implies Lipschitz}
\end{itemize}
\end{dependencies}

\begin{figure}[H]
\centering
\resizebox{0.78\linewidth}{!}{%
\begin{tikzpicture}[atlas,scale=1.05]
  \def\f{1.8+0.45*sin(70*\x)}
  \glowcurve{color=atlasblue, domain=0.3:4.7, axis=0, top=3.4, expr={\f}}
  \draw[atlas axes] (0,0)--(5.1,0);
  \draw[atlas axes] (0,0)--(0,3.4);
  \def\p{2.0}
  \pgfmathsetmacro\fp{1.8+0.45*sin(70*\p)}
  \pgfmathsetmacro\mp{0.45*70*0.0174533*cos(70*\p)}
  \def\L{0.7}
  \draw[atlas probe,atlasconegreen,dash pattern=on 3pt off 2pt]
    plot[domain=0.45:3.55,samples=80,smooth]
    (\x,{\fp+\mp*(\x-\p)+0.5*\L*(\x-\p)*(\x-\p)});
  \draw[atlas probe,atlasconegreen,dash pattern=on 3pt off 2pt]
    plot[domain=0.45:3.55,samples=80,smooth]
    (\x,{\fp+\mp*(\x-\p)-0.5*\L*(\x-\p)*(\x-\p)});
  \draw[atlas probe,atlasred]
    ({\p-1.4},{\fp-\mp*1.4})--({\p+1.4},{\fp+\mp*1.4});
  \node[atlas dot] at (\p,\fp){};
  \draw[dropline] (\p,\fp)--(\p,0);
  \node[atlas label,below] at (\p,-0.05){\footnotesize$p$};
  \node[atlas label,above] at (2.6,3.15){\scriptsize$|f'(x)-f'(y)|\le L\,|x-y|$};
  \node[atlas label,atlasconegreen!60!atlasink,above left] at (1.0,2.55)
    {\scriptsize$f(p)+f'(p)(x-p)+\tfrac{L}{2}(x-p)^2$};
  \node[atlas label,atlasconegreen!60!atlasink,below right=2pt] at (3.0,0.55)
    {\scriptsize$-\tfrac{L}{2}(x-p)^2$};
  \node[atlas label,atlasred!70!atlasink,below left] at (\p-0.5,{\fp-\mp*0.5})
    {\scriptsize tangent at $p$};
  \node[atlas label,right] at (3.9,2.35){\footnotesize$y=f(x)$};
\end{tikzpicture}%
}
\caption{\(C^{1,1}\) control gives quadratic upper and lower bounds around a tangent line.}
\label{fig:c11-descent-bounds}
\end{figure}


\begin{corollarybox}{Corollary (Bounded Second Derivative Implies \(C^{1,1}\))}
\begin{corollary}[Bounded Second Derivative Implies \(C^{1,1}\)]
\label{cor:bounded-second-derivative-implies-c11}
Let \(I\subseteq\mathbb{R}\) be an interval and let \(f\in C^2(I)\). If
\(|f''(x)|\leq M\) for all \(x\in I\), then \(f\in C^{1,1}(I)\), and \(f'\)
is Lipschitz with constant \(M\).
\hyperref[prf:bounded-second-derivative-implies-c11]{Go to proof.}
\end{corollary}
\end{corollarybox}

\begin{remark*}[Standard quantified statement]
\[
  \bigl(f\in C^2(I)\land \forall x\in I:\ |f''(x)|\leq M\bigr)
  \Rightarrow
  f\in C^{1,1}(I)\land \operatorname{IsLipschitz}(f',I,M).
\]
\end{remark*}

\begin{remark*}[Interpretation]
This is \hyperref[cor:derivative-bound-implies-lipschitz]{Derivative Bound
Implies Lipschitz} applied one level up, to \(f'\) in place of \(f\). The
converse requires almost-everywhere language and is therefore deferred.


A Lipschitz derivative is differentiable at most points, but making that
sentence precise requires measure zero. In one variable this passes through
bounded variation; in several variables it becomes Rademacher's theorem.
Those results belong to the integration and measure-theoretic development,
not to this chapter.
\end{remark*}

\begin{dependencies}
\begin{itemize}
  \item \hyperref[cor:derivative-bound-implies-lipschitz]{Corollary: Derivative Bound Implies Lipschitz}
  \item \hyperref[def:class-ck]{Definition: Class \(C^k\)}
  \item \hyperref[def:class-c11]{Definition: Lipschitz-Derivative Class \(C^{1,1}\)}
\end{itemize}
\end{dependencies}
